\documentclass{article}

% Language setting
% Replace `english' with e.g. `spanish' to change the document language
\usepackage[french]{babel}
\usepackage[T1]{fontenc}
\usepackage{enumitem}
\setlist[itemize]{label=\textbullet}

% Set page size and margins
% Replace `letterpaper' with `a4paper' for UK/EU standard size
\usepackage[letterpaper,top=2cm,bottom=2cm,left=3cm,right=3cm,marginparwidth=1.75cm]{geometry}

% Useful packages
\usepackage{amsmath}
\usepackage{graphicx}
\usepackage[colorlinks=true, allcolors=blue]{hyperref}

\hfuzz=\maxdimen %suppress overfull hbox

\title{Recherche bibliographique}
\author{}
\date{}

\begin{document}
\maketitle
\tableofcontents
% créer des sections à mesure que je découvre des articles.

\newpage
\section{Santé humaine}
    \subsection{Viral oncogenesis in cancer: from mechanisms to therapeutics}
    \noindent
    \url{https://doi.org/10.1038/s41392-025-02197-9}
        \subsubsection{Abstract}

\begin{itemize}
    \item Virus d'Epstein–Barr (1964).
    \item Premier virus oncogène identifié.
    \item Hépatite B et papillomavirus : virus oncogènes. Certains cancers sont associés au VIH (ex. : sarcome de Kaposi).
    \item Transformation cellulaire et altération du métabolisme cellulaire.
    \item Comprendre les mécanismes oncogènes pour mieux traiter.
\end{itemize}

\subsubsection{Corps de l'article}
\begin{itemize}
    \item En 1964, Epstein et Barr ont identifié l'EBV chez des patients atteints de lymphome de Burkitt. L'EBV a ensuite été associé à la mononucléose infectieuse et au carcinome nasopharyngé.
    \item En 1953, Rowe et al. ont décrit les adénovirus.
    \item EBV : grand génome ADN double brin.
    \item HBV : petit génome ADN partiellement double brin.
    \item HCV : génome ARN simple brin.
    \item HIV : rétrovirus.
    \item Remarque : les virus à ARN possèdent un génome en ARN (simple ou double brin), tandis que les rétrovirus ont un génome ARN simple mais passent par un intermédiaire ADN au cours de leur cycle.
    \item Ceci suggère que l'infection n'est qu'un facteur de l'oncogenèse, qui est un processus multifactoriel et complexe.
    \item \textbf{HBV :} virus ADN hépatotrope partiellement double brin ; l'enveloppe encapsule un nucléocapside contenant de l'ADN circulaire relaxé (rcDNA). Le rcDNA est transloqué au noyau et converti en ADN circulaire covalently fermé (cccDNA). Le cccDNA sert de matrice pour la transcription du pré‑ARN génomique (pgRNA) et des ARNm ; la rétrotranscription du pgRNA produit l'ADN HBV et les protéines virales. L'assemblage de l'enveloppe se fait dans des corps multivésiculaires (MVB) pour produire et sécréter les particules virales.
    \item \underline{Virus ADN :}
    \item \textbf{HPV :} virus non enveloppé à ADN double brin circulaire ; certains types (HPV16, HPV18) sont oncogènes, d'autres (HPV6, HPV11) le sont moins. Environ 200 types. Les gènes E1 et E2 interviennent dans la réplication, L1 et L2 dans l'empaquetage ; E4, E5, E6 et E7 modulent le cycle cellulaire, l'évasion immunitaire et la libération virale. Le HPV infecte principalement les cellules basales de l'épithélium en se liant à la matrice extracellulaire ou à la surface cellulaire ; après internalisation la capside est désassemblée dans les endosomes et L2 guide l'ADN viral vers le noyau. Une fois l'infection établie, le cycle viral est régulé pour maintenir l'homéostasie protéique hôte‑virus.
    \item \textbf{EBV :} herpesvirus gamma, capside + enveloppe, ADN linéaire double brin ; cibles : lymphocytes et cellules épithéliales. L'entrée dans les cellules B débute par la liaison de gp350 au récepteur CD21 ; gp42 interagit ensuite avec le CMH‑II, permettant l'action du complexe gH/gL qui active la protéine de fusion gB et induit la fusion membranaire. L'infection peut être latente ou lytique : quatre types de latence (III, II, I, 0) correspondant à une restriction progressive de l'expression génique pour échapper à la surveillance immunitaire. L'EBV peut établir une persistance dans les cellules B mémoire (latence 0) et se réactiver périodiquement, contribuant à la réplication virale et à l'oncogenèse. Les types de latence sont associés à différents cancers.
    \item \textbf{KSHV :} herpesvirus humain gamma, ADN linéaire double brin avec unités répétées terminales riches en GC ; infecte les cellules B, les fibroblastes, les cellules épithéliales et endothéliales. Présent en états latent et lytique ; la protéine LANA maintient la latence. Divers stimuli peuvent provoquer la sortie de latence ; RTA régule la transition vers la réplication lytique.
    \item \begin{itemize}
        \item \texttt{Épithélium} : couche cellulaire recouvrant les surfaces externes (peau) et internes (parois d'organes, cavités) ; fonctions : protection, absorption, sécrétion, transport.
        \item \texttt{Endothélium} : épithélium simple pavimenteux spécialisé tapissant l'intérieur des vaisseaux sanguins et lymphatiques et la cavité cardiaque ; fonctions : perméabilité, hémostase, régulation du tonus vasculaire, trafic leucocytaire.
    \end{itemize}
    \item \underline{Virus ARN :}
    \item \textbf{HCV :} virus ARN simple brin sens positif hépatotrope, enveloppé, avec nucléocapside. Les glycoprotéines d'enveloppe E1 (fusion) et E2 (liaison au récepteur) jouent des rôles clés. L'ARN polymérase virale sans correction d'erreur conduit à une grande variabilité génétique ; les anticorps ciblent principalement E2/NS1. La lyse des cellules infectées est médiée par la perforine et la granzyme B des LT cytotoxiques et des NK. La majorité des infections évoluent vers une infection chronique.
    \item \textbf{HIV :} virus ARN simple brin infectant principalement les lymphocytes CD4+, les macrophages et les cellules dendritiques. Composants : polyprotéine Gag et enveloppe. La liaison de gp120 aux corécepteurs CCR5/CXCR4 initie l'entrée ; la transcriptase inverse produit de l'ADN viral qui s'intègre au génome hôte, formant des réservoirs latents. L'absence de correction par la RT favorise les mutations. Gag coordonne l'assemblage et la libération des particules virales. Phases d'infection : aiguë, asymptomatique, sida.
    \item \underline{Conclusion :}
    \item Divers oncovirus pénètrent les cellules par liaison à des récepteurs spécifiques via des protéines structurales. Les virus ADN (HBV, HPV, EBV, KSHV, MCV) peuvent intégrer directement leur génome ; les rétrovirus (HIV, HTLV‑1) nécessitent une transcription inverse avant intégration.
    \item Après intégration, un oncovirus peut activer des proto‑oncogènes ou inactiver des gènes suppresseurs de tumeur, déclenchant la carcinogenèse. Bien que le HCV n'intègre pas son génome, l'inflammation chronique augmente le risque de mutations et de cancer.
    \item Mécanismes de carcinogenèse : instabilité génomique, dérégulation des oncogènes et suppresseurs, perturbation du cycle cellulaire, immunosuppression, reprogrammation métabolique et inflammation chronique.
    \item Les virus ADN entraînent souvent la transformation via des oncoprotéines virales qui altèrent les voies de signalisation et le cycle cellulaire.
    \item Les virus ARN favorisent l'oncogenèse surtout par infection chronique, inflammation persistante et immunosuppression, facilitant l'accumulation de mutations ; le VIH agit principalement indirectement.
    \item Les thérapies ciblées contre les tumeurs virales : vaccins, thérapies géniques ciblées modifiant des vecteurs viraux, chimiothérapie.
    \item Les virus oncogènes (ex. HPV) présentent une grande variabilité génétique, compliquant l'identification de cibles thérapeutiques stables et réduisant l'efficacité des traitements face aux mutations.
    \item Les cancers associés aux virus développent souvent des mécanismes d'évasion immunitaire complexes (ex. : HCC induit par HBV diminuant l'expression du CMH).
    \item Les vaccins resteront une avancée cruciale.
    \item Les immunothérapies et thérapies ciblées actuelles montrent des bénéfices limités ; une meilleure compréhension du système immunitaire pourrait améliorer les résultats via les inhibiteurs de points de contrôle et les lymphocytes T génétiquement modifiés.
\end{itemize}

\newpage
\section{Biologie animale/Bactériologie/Virologie}

\newpage
\section{Biologie végétale/Mycologie/Ecologie}

\newpage
\section{Informatique}


\bibliographystyle{alpha}
\bibliography{sample}

\end{document}